%%%%%%%%%%%%%%%%%%%%%%%%%%%%%%%%%%%%%%%%%%%%%%%%%%%%%%%%%%%%%%%%%%%%%%%%%%%%%%%
% APPENDIX BA: FORMAL-SEGMENTATION — EMERGENT SEGMENT GENESIS
%%%%%%%%%%%%%%%%%%%%%%%%%%%%%%%%%%%%%%%%%%%%%%%%%%%%%%%%%%%%%%%%%%%%%%%%%%%%%%%
%
% CATEGORY: FORMAL (Mathematical Formalization of Segmentation Genesis)
% Module: BCM2_09_SEG7257 (Segmentlogik – Emergente Segmentgenese)
% Category: FORMAL
% Version: 2.0 (January 2026)
% Status: Final
% Language: English
%
% FUNDAMENTAL QUESTION: "How do behavioral segments emerge as natural clusters
% from heterogeneity in utility preferences, awareness, willingness, and context?"
%
% PURPOSE:
% Complete formal axiomatization of segment emergence mechanisms, showing that
% populations naturally self-organize into behavioral clusters based on the
% underlying distributions of utility weights (ω), awareness (A), willingness (W),
% and context (Ψ). Segment types and counts are NOT fixed a priori but emerge
% dynamically from these distributions. This formalization is universally applicable
% across all behavior-change domains and contexts.
%
% CHAPTER LINKAGE:
%   Primary: Chapter 14: Behavioral Change Segments (BCS)
%   Secondary: Chapters 11-13 (Awareness → Willingness → Journey)
%
%%%%%%%%%%%%%%%%%%%%%%%%%%%%%%%%%%%%%%%%%%%%%%%%%%%%%%%%%%%%%%%%%%%%%%%%%%%%%%%

\chapter{FORMAL-SEGMENTATION: Emergent Segment Genesis}
\label{app:formal-segmentation}
\label{app:ba-segments}

%------------------------------------------------------------------------------
% HEADER BLOCK
%------------------------------------------------------------------------------

\begin{tcolorbox}[colback=gray!5!white, colframe=gray!75!black,
    title=Appendix PAP1: Header Block]

\noindent
\textbf{Code:} BA \quad \textbf{Category:} FORMAL-SEGMENTATION \\
\textbf{Module:} BCM2\_09\_SEG7257 \quad \textbf{Version:} 2.0 (January 2026) \\
\textbf{Status:} Final \quad \textbf{Language:} English

\vspace{0.5em}

\noindent
\textbf{Description:}
Formal axiomatization of behavioral segment emergence from population heterogeneity. Shows that segments are natural clusters emerging from distributions of utility preferences, awareness, and willingness—not fixed archetypes. Universally applicable across all behavior-change domains.

\end{tcolorbox}

%------------------------------------------------------------------------------
% CHAPTER LINKAGE BOX
%------------------------------------------------------------------------------

\begin{tcolorbox}[colback=purple!5!white, colframe=purple!75!black,
    title=Chapter Linkage]

\textbf{Primary Chapter:}
\begin{itemize}[nosep]
\item Chapter 14: Behavioral Change Segments $\leftrightarrow$ Sections \ref{sec:ba-genesis}--\ref{sec:ba-context-dependence}
\end{itemize}

\textbf{Secondary Chapters:}
\begin{itemize}[nosep]
\item Chapter 11: Awareness --- provides $A(t)$ distributions
\item Chapter 12: Willingness --- provides $W(t)$ framework
\item Chapter 13: Behavioral Change Journey --- provides BCJ-Stage dynamics
\end{itemize}

\textbf{Reading Guide:}
\begin{itemize}[nosep]
\item For conceptual overview: Read Chapter 14 first
\item For formal axiomatization: This appendix provides complete specification
\item For implementation: See Chapter 14 Section on Clustering Approach
\end{itemize}

\end{tcolorbox}

%------------------------------------------------------------------------------
% ABSTRACT
%------------------------------------------------------------------------------

\begin{abstract}
This appendix provides formal axiomatization of behavioral segment emergence.
Rather than positing fixed archetypal segments, we show how segments naturally
emerge as clusters in the multidimensional space of utility weights ($\omega_i$),
awareness ($A_i$), willingness ($W_i$), and context ($\Psi$). The fundamental
mechanism is universal: heterogeneity in individual preferences and decision
processes creates natural clustering. However, the number, type, and size of
clusters are context-dependent and domain-specific. This framework explains why
some domains show 3-4 natural segments while others show 6-7. The formalization
is based on 6 core axioms governing segment emergence, plus 5 propositions about
cluster stability, transitions, and context modulation.
\end{abstract}

%------------------------------------------------------------------------------
% QUICK REFERENCE BOX
%------------------------------------------------------------------------------

\begin{tcolorbox}[colback=blue!5!white, colframe=blue!75!black,
    title=Quick Reference: Segmentation Genesis]

\textbf{Core Question:} How do natural behavioral segments emerge from
heterogeneous populations?

\textbf{Key Objects:}
Population heterogeneity space:
\begin{equation}
\mathcal{H} = \{\omega_i, A_i, W_i, \Psi | i \in [1,N]\}
\end{equation}

\textbf{Central Mechanism:}
\begin{equation}
\text{Segments} = \text{Clusters}(\mathcal{H}, \text{distance metric}, k)
\end{equation}

where $k$ (segment count) is data-driven, not fixed.

\textbf{Universal Axioms:} 6 axioms govern emergence (not 11)

\textbf{Related Appendices:} AW (STAGE), AV (READY), AU (AWARE), C (WHAT)

\end{tcolorbox}

%%%%%%%%%%%%%%%%%%%%%%%%%%%%%%%%%%%%%%%%%%%%%%%%%%%%%%%%%%%%%%%%%%%%%%%%%%%%%%%
% -----------------------------------------------------------------------------
% APPENDIX SCOPE BOX
% -----------------------------------------------------------------------------

\begin{tcolorbox}[
    colback=orange!5!white,
    colframe=orange!75!black,
    title={\textbf{Appendix Scope: Ziel / In-Scope / Out-of-Scope / Constraints}},
    fonttitle=\bfseries
]

\textbf{Ziel (Objective):}
\begin{itemize}[nosep]
    \item Given a population with heterogeneous utility weights (ω), awareness (A),
willingness thresholds (W), and context (Ψ), what formal principles govern how
behavioral segments emerge as natural clusters? And how do these principles
apply universally across all domains, while producing domain-specific segment
structures?
\end{itemize}

\textbf{In-Scope:}
\begin{itemize}[nosep]
    \item The Fundamental Problem
    \item Axioms of Segment Emergence
    \item Why 5 Segments in Some Domains: A Special Case
    \item Integration with CORE Framework
    \item Worked Example: Environmental Behavior in Three Contexts
\end{itemize}

\textbf{Out-of-Scope (delegiert an andere Appendices):}
\begin{itemize}[nosep]
    \item [TODO: Define delegations to other appendices]
\end{itemize}

\textbf{Constraints (Anwendungsgrenzen):}
\begin{itemize}[nosep]
    \item [TODO: Define when this appendix does NOT apply]
    \item [TODO: Define required prerequisites]
    \item [TODO: Define known limitations]
\end{itemize}

\textbf{Lieferobjekte (Deliverables):}
\begin{itemize}[nosep]
    \item 4 Table(s)
    \item 35 Equation(s)
    \item Worked Example(s)
\end{itemize}

\end{tcolorbox}


\section{The Fundamental Problem}
\label{sec:ba-fundamental}
%%%%%%%%%%%%%%%%%%%%%%%%%%%%%%%%%%%%%%%%%%%%%%%%%%%%%%%%%%%%%%%%%%%%%%%%%%%%%%%

\subsection{Why Universal Segmentation Theory Fails}

Classical behavior-change models assume homogeneity: all decision-makers have
the same utility function, awareness capacity, and action thresholds. Empirically,
this fails. Real populations show massive heterogeneity.

But a deeper problem exists: **we cannot enumerate all possible segment types
in advance**. Different domains produce different natural clusters:

\begin{itemize}
\item \textbf{Environmental behavior}: ~5 natural segments (Rogers diffusion curve)
\item \textbf{Healthcare adoption}: ~3-4 segments (medical necessity dominates)
\item \textbf{Financial decisions}: ~6-7 segments (risk profiles vary widely)
\item \textbf{Organizational change}: ~4-5 segments (role-driven variation)
\end{itemize}

The number varies! Why?

\subsection{The True Universal Principle}

Not: ``There are these 5 segments.''

But: **``From heterogeneous populations, natural clusters emerge. The structure
of these clusters depends on the underlying distributions of preferences and
decision processes.''**

This is universal. The details are domain-specific.

\begin{tcolorbox}[colback=red!5!white, colframe=red!75!black,
    title=The Fundamental Question]
\textbf{Given a population with heterogeneous utility weights (ω), awareness (A),
willingness thresholds (W), and context (Ψ), what formal principles govern how
behavioral segments emerge as natural clusters? And how do these principles
apply universally across all domains, while producing domain-specific segment
structures?}
\end{tcolorbox}

%%%%%%%%%%%%%%%%%%%%%%%%%%%%%%%%%%%%%%%%%%%%%%%%%%%%%%%%%%%%%%%%%%%%%%%%%%%%%%%
\section{Axioms of Segment Emergence}
\label{sec:ba-genesis}
%%%%%%%%%%%%%%%%%%%%%%%%%%%%%%%%%%%%%%%%%%%%%%%%%%%%%%%%%%%%%%%%%%%%%%%%%%%%%%%

This section formalizes 6 universal axioms of segment emergence (not domain-specific).

%----------- AXIOM 1 -----------
\begin{tcolorbox}[colback=gray!5!white, colframe=gray!75!black,
    title=Axiom SEG-EMERGE-1: Universal Population Heterogeneity]

\textbf{Statement:}
Every population exhibits heterogeneity in utility preferences, awareness levels,
and willingness thresholds. This heterogeneity is \textit{not an exception but
the universal norm}.

\textbf{Formal Definition:}
For any population $N$, define individual-level parameters:
\begin{equation}
\vec{\omega}_i = (\omega^F_i, \omega^E_i, \omega^P_i, \omega^S_i, \omega^D_i, \omega^E_i)
\quad \text{utility weights for individual } i
\end{equation}
\begin{equation}
A_i(t) \in [0,1] \quad \text{awareness trajectory}
\end{equation}
\begin{equation}
W_i(t) \quad \text{willingness function}
\end{equation}

Heterogeneity exists if the variance across individuals is non-negligible:
\begin{equation}
\text{Var}_i(\omega^F_i) > \varepsilon, \quad \text{Var}_i(A_i) > \varepsilon, \quad
\text{Var}_i(W_i) > \varepsilon
\end{equation}

for threshold $\varepsilon > 0$ (typically, empirical variance is large).

\textbf{Interpretation:}
No behavioral model can explain behavior by averaging parameters across the
population. Individual differences matter.

\textbf{Implication:}
Population-level behavior emerges from aggregating heterogeneous individuals,
not from applying a representative-agent model.

\end{tcolorbox}

\vspace{1em}

%----------- AXIOM 2 -----------
\begin{tcolorbox}[colback=gray!5!white, colframe=gray!75!black,
    title=Axiom SEG-EMERGE-2: Natural Clustering in Preference Space]

\textbf{Statement:}
Heterogeneous populations do not distribute uniformly in preference space. Instead,
they naturally cluster into groups of individuals with similar preferences. These
clusters are the foundation of behavioral segments.

\textbf{Formal Definition:}
Define the \textit{behavioral heterogeneity space}:
\begin{equation}
\mathcal{H} = \left\{ (\vec{\omega}_i, A_i, W_i) \in \mathbb{R}^{6 + 1 + 1} \right\}
\end{equation}

(6 utility dimensions, 1 awareness level, 1 willingness parameter per individual)

A \textit{behavioral segment} is a cluster $C_k \subset \mathcal{H}$ such that:
\begin{equation}
C_k = \{ i \in [1,N] : d((\vec{\omega}_i, A_i, W_i), \mu_k) < \theta_k \}
\end{equation}

where:
\begin{itemize}[nosep]
\item $d(\cdot, \cdot)$ is a distance metric in preference space
\item $\mu_k$ is the cluster center (prototype)
\item $\theta_k$ is the cluster radius (intra-cluster cohesion threshold)
\end{itemize}

Clusters are natural if they minimize within-cluster variance and maximize
between-cluster variance (standard clustering objective).

\textbf{Interpretation:}
Segments are not imposed from theory but emerge from the data. Individuals with
similar preference structures naturally group.

\textbf{Implication:}
Segmentation is fundamentally a \textit{data-driven} (not a priori) operation.
The number of clusters $k$ depends on the data distribution, not on theoretical
assumptions.

\end{tcolorbox}

\vspace{1em}

%----------- AXIOM 3 -----------
\begin{tcolorbox}[colback=gray!5!white, colframe=gray!75!black,
    title=Axiom SEG-EMERGE-3: Context Modulates Cluster Structure}

\textbf{Statement:}
The cluster structure (number of clusters $k$, cluster sizes $|C_k|$, cluster
prototypes $\mu_k$) is not invariant across contexts. Context $\Psi$ directly
modulates the heterogeneity distribution, thus changing the clustering structure.

\textbf{Formal Definition:}
The clustering is conditional on context:
\begin{equation}
\mathcal{C} = \text{Clusters}(\mathcal{H} | \Psi)
\end{equation}

Different contexts produce different heterogeneity distributions:
\begin{equation}
f(\mathcal{H} | \Psi_1) \neq f(\mathcal{H} | \Psi_2)
\quad \text{for } \Psi_1 \neq \Psi_2
\end{equation}

Examples:
\begin{itemize}[nosep]
\item High $\Psi_{\text{Uncertainty}}$ (novel technology): Awareness heterogeneity increases,
  creating more risk-averse vs. adventurous clusters
\item High $\Psi_{\text{Social}}$ (strong peer effects): Social-preference heterogeneity
  increases, creating more social-conformist clusters
\item High $\Psi_{\text{Economic}}$ (expensive behavior): Financial-concern heterogeneity
  increases, creating more pragmatist clusters
\end{itemize}

\textbf{Interpretation:}
Segments are not static categories but dynamic structures shaped by context.
A person can be in Segment A in Context 1 and Segment B in Context 2 if the
context changes their relative preference structure.

\textbf{Implication:}
Effective segmentation analysis must be domain- and context-specific. Generic
segment models fail.

\end{tcolorbox}

\vspace{1em}

%----------- AXIOM 4 -----------
\begin{tcolorbox}[colback=gray!5!white, colframe=gray!75!black,
    title=Axiom SEG-EMERGE-4: Cluster-Dependent Behavior Change Dynamics}

\textbf{Statement:}
Individuals within the same cluster have similar response patterns to awareness,
willingness interventions, and social influence. Between clusters, response
patterns differ systematically.

\textbf{Formal Definition:}
Define intervention response as elasticity:
\begin{equation}
\epsilon_i^{I} = \frac{\partial W_i}{\partial I}
\quad \text{(willingness elasticity to intervention I)}
\end{equation}

Cluster homogeneity implies:
\begin{equation}
\text{Var}_{i \in C_k}(\epsilon_i^{I}) < \text{Var}_{i \in [1,N]}(\epsilon_i^{I})
\end{equation}

That is: within-cluster variance in response patterns is lower than population
variance. This makes segments useful for targeted intervention design.

Cluster-to-cluster differences imply:
\begin{equation}
E[\epsilon_i^{I} | i \in C_k] \neq E[\epsilon_i^{I} | i \in C_k']
\quad \text{for } k \neq k'
\end{equation}

\textbf{Interpretation:}
If Axiom 4 did not hold (clusters had random response patterns), segmentation
would be useless for intervention design. The practical power of segmentation
derives from cluster-level behavioral homogeneity.

\textbf{Implication:}
Validation of a segmentation scheme requires empirical verification: Do clusters
show statistically significant differences in intervention response?

\end{tcolorbox}

\vspace{1em}

%----------- AXIOM 5 -----------
\begin{tcolorbox}[colback=gray!5!white, colframe=gray!75!black,
    title=Axiom SEG-EMERGE-5: Context-Contingent Cluster Transitions}

\textbf{Statement:}
Individuals can transition between clusters as their awareness, willingness,
and preferences evolve. Transition probabilities are shaped by cluster structure
and context.

\textbf{Formal Definition:}
Define the transition probability from cluster $C_i$ to cluster $C_j$:
\begin{equation}
P(C_i \to C_j | t) = f\left( d(\vec{\omega}_i(t), \mu_j), A_i(t), W_i(t), \Psi(t) \right)
\end{equation}

Key properties:
\begin{itemize}[nosep]
\item $P(C_i \to C_j)$ increases as individual $i$'s preferences $\vec{\omega}_i(t)$ move
  closer to cluster $j$'s prototype $\mu_j$
\item Transitions are more likely when awareness increases (individual can perceive
  alternative preference clusters)
\item Transitions are more likely in stable, low-uncertainty contexts (individual
  can safely experiment with new preferences)
\end{itemize}

\textbf{Interpretation:}
Clusters are stable attractors but not permanent prisons. Behavioral change can
move individuals between clusters.

\textbf{Implication:}
Interventions can trigger inter-cluster transitions. This is more efficient than
trying to change behavior \textit{within} a cluster.

\end{tcolorbox}

\vspace{1em}

%----------- AXIOM 6 -----------
\begin{tcolorbox}[colback=gray!5!white, colframe=gray!75!black,
    title=Axiom SEG-EMERGE-6: Cluster Count Determination}

\textbf{Statement:}
The number of natural clusters $k$ is not fixed a priori but determined by the
underlying heterogeneity distribution. Different distributions support different
values of $k$. Standard clustering algorithms (k-means, hierarchical, spectral)
can identify $k$ empirically.

\textbf{Formal Definition:}
Optimal $k$ minimizes a cluster-quality criterion such as silhouette score,
within-cluster-sum-of-squares (WSS), or Bayesian Information Criterion (BIC):
\begin{equation}
k^* = \arg\min_k \left[ WSS(k) + \lambda \cdot k \right]
\quad \text{or equivalently}
\quad
k^* = \arg\max_k \text{Silhouette}(k)
\end{equation}

The optimal $k$ is typically in range $[2, 8]$ for most behavior-change domains,
but can fall outside this range in domain-specific contexts.

\textbf{Interpretation:}
No theory mandates ``5 segments are universal.'' Instead, empirical analysis of
each population/context determines the natural number of clusters. This is why
different behavior-change domains show different segment structures.

\textbf{Implication:}
Segmentation analysis must include a \textit{k-determination step} (typically
via elbow method, silhouette analysis, or gap statistic) rather than imposing
a fixed count.

\end{tcolorbox}

\vspace{1em}

%----------- AXIOM 7 -----------
\begin{tcolorbox}[colback=gray!5!white, colframe=gray!75!black,
    title=Axiom SEG-EMERGE-7: Natural Distribution of Cluster Sizes Under Heterogeneity]

\textbf{Statement:}
The distribution of cluster sizes $|C_k|$ (number of individuals in cluster $k$) is
determined by the underlying distribution of heterogeneity in preference space.
Under weak regularity conditions, cluster sizes follow predictable patterns that
depend on the shape of the heterogeneity distribution $f(\mathcal{H} | \Psi)$.

\textbf{Formal Definition:}
Let the heterogeneity distribution be approximately multivariate normal:
\begin{equation}
f(\mathcal{H} | \Psi) \approx \mathcal{N}(\mu_{\mathcal{H}}, \Sigma_{\mathcal{H}})
\end{equation}

Then, under standard clustering of this normal distribution, cluster sizes follow:
\begin{equation}
|C_k| \sim \text{Binomial-like distribution with } E[|C_k|] \approx \frac{N}{k^*}
\end{equation}

The coefficient of variation in cluster sizes is:
\begin{equation}
CV(|C_k|) = \frac{\sigma_{n_k}}{\bar{n}_k} \approx 0.3 \text{ to } 0.5
\quad \text{(typically moderate variability)}
\end{equation}

For the special case of **adoption diffusion with social learning** (Rogers structure),
cluster sizes approximately follow a **Beta distribution** with parameters aligned to:
\begin{equation}
|C_k| \propto \text{Beta}(\alpha, \beta) \quad \text{with } \alpha \approx 2, \beta \approx 3
\end{equation}

This produces the characteristic distribution: $\{2.5\%, 13.5\%, 34\%, 34\%, 16\%\}$.

\textbf{Interpretation:}
Cluster sizes are not arbitrary. They emerge from the geometry of the heterogeneity
distribution. Approximately normal heterogeneity produces relatively balanced cluster
sizes (CV ≈ 0.3-0.5). Highly skewed or bimodal heterogeneity produces very different
patterns (e.g., two dominant clusters for polarized populations).

\textbf{Implication:}
If cluster sizes deviate radically from expected patterns under normality (e.g., one
cluster has 70% and others have <10%), this signals either:
\begin{enumerate}[nosep]
\item The underlying heterogeneity distribution is non-normal (bimodal, heavily skewed)
\item The clustering algorithm is suboptimal for this data structure
\item The context has changed, producing a different heterogeneity distribution
\end{enumerate}

Practitioners can use cluster size distributions as **diagnostic indicators** of
whether segmentation is capturing natural structure or artifacts.

\end{tcolorbox}

%----------- AXIOM 8 -----------
\begin{tcolorbox}[colback=red!5!white, colframe=red!75!black,
    title=Axiom SEG-MATCH-8: Intervention Leverage Principle]

\textbf{Statement:}
Behavioral clusters have \textit{identifiable leverage points}---specific utility dimensions
and awareness barriers that interventions can strategically address. The effectiveness of
an intervention on a cluster depends on the degree to which the intervention targets
that cluster's dominant constraints and leverageable dimensions.

\textbf{Formal Definition:}
For each cluster $C_k$ in a given context, define the \textbf{leverage profile} as:
\begin{equation}
\mathcal{L}_k = (\ell_k^d)_{d \in \{F,E,P,S,D,E\}} \quad \text{where} \quad \ell_k^d = \frac{\text{Var}(\omega_i^d | i \in C_k)}{\sum_{d'} \text{Var}(\omega_i^{d'} | i \in C_k)}
\end{equation}

The leverage profile identifies which utility dimensions show the \textit{most variability within the cluster}:
- High $\ell_k^d$ means cluster members differ most on dimension $d$ (high leverage for interventions targeting $d$)
- Low $\ell_k^d$ means cluster members are aligned on dimension $d$ (low leverage, targeting this dimension wastes intervention effort)

Additionally, define the \textbf{constraint profile} as:
\begin{equation}
\mathcal{C}_k = (\gamma_k^b)_{b \in \{\text{barriers}\}} \quad \text{where} \quad \gamma_k^b = \Pr(\text{barrier } b \text{ blocks } i \in C_k)
\end{equation}

Common barriers include: $(b) =$ \{low awareness, low willingness, high implementation costs, social disapproval, habit inertia, etc.\}

\textbf{Intervention Leverage Matching:}
An intervention $I$ can be described by its \textbf{design profile}:
\begin{equation}
\mathcal{D}_I = (\delta_I^d, \beta_I^b) \quad \text{where}
\end{equation}
\begin{itemize}[nosep]
\item $\delta_I^d \in [0,1]$ = intensity of dimension $d$ in intervention (which motivations does it appeal to?)
\item $\beta_I^b \in [0,1]$ = effectiveness of intervention at addressing barrier $b$
\end{itemize}

For an intervention to have leverage on cluster $C_k$, it must target the cluster's actual constraints:
\begin{equation}
\text{Leverage}(I, C_k) = \sum_d (\ell_k^d \cdot \delta_I^d) + \sum_b (\gamma_k^b \cdot \beta_I^b)
\end{equation}

\textbf{Interpretation:}
High leverage means the intervention is well-matched to the cluster's profile:
- It appeals to dimensions on which cluster members vary most
- It addresses the barriers actually blocking the cluster
- Expected impact is proportional to leverage

Low leverage means poor matching:
- Intervention appeals to dimensions on which cluster is already aligned
- Intervention addresses barriers the cluster doesn't face
- Expected impact approaches zero

\textbf{Critical Implication for Your Question:}
\textbf{If segmented and non-segmented groups perform equally, compute Leverage$(I, C_k)$ for each cluster.}
If average leverage is low, the problem is \textbf{poor intervention matching}, not bad segmentation.
If average leverage is high but outcomes still equal, then segmentation itself may be weak (Axioms 1-7 fail).

\textbf{Example:}
\begin{itemize}[nosep]
\item Cluster with high $\omega_{\text{Identity}}$: Will respond to identity-framed messages ($\delta_I^{S} \approx 1.0$)
\item Cluster with high $\omega_{\text{Financial}}$: Will respond to ROI/cost-benefit messages ($\delta_I^{E} \approx 1.0$)
\item Mismatch: Identity-driven cluster + Financial-only intervention = $\text{Leverage}(I, C_k) \approx 0.1$ (poor match)
\end{itemize}

\end{tcolorbox}

\vspace{1em}

%----------- AXIOM 9 -----------
\begin{tcolorbox}[colback=blue!5!white, colframe=blue!75!black,
    title=Axiom SEG-MATCH-9: Specificity-Matching Principle]

\textbf{Statement:}
The expected effect size of an intervention on a cluster is a function of the match between
the intervention's design and the cluster's leverage profile. Interventions matched to cluster
characteristics produce significantly larger effects than generic (non-segmented) interventions.
The magnitude of this effect is predictable and measurable.

\textbf{Formal Definition:}
Define the \textbf{expected outcome change} for cluster $k$ under intervention $I$:
\begin{equation}
E[\Delta O_k^I] = E[\Delta O_{\text{baseline}}] + \beta^{\text{match}} \cdot \text{Leverage}(I, C_k) + \epsilon_k
\end{equation}

Where:
\begin{itemize}[nosep]
\item $E[\Delta O_{\text{baseline}}]$ = baseline effect when no segmentation (same intervention for all)
\item $\beta^{\text{match}}$ = effect size multiplier per unit of leverage (estimated empirically)
\item $\text{Leverage}(I, C_k)$ = intervention-cluster match (from Axiom 8)
\item $\epsilon_k$ = cluster-specific residual (measurement error, implementation fidelity, etc.)
\end{itemize}

In standardized effect size terms:
\begin{equation}
\text{Cohen's } d_k^I = d_{\text{baseline}} + \Delta d_{\text{match}} \cdot \frac{\text{Leverage}(I, C_k)}{\max_k \text{Leverage}(I, C_k)}
\end{equation}

\textbf{Prediction Framework:}
Using this axiom, we can predict effect patterns:

\begin{table}[h]
\centering
\small
\begin{tabular}{lll}
\toprule
\textbf{Leverage Match} & \textbf{Predicted Effect} & \textbf{Expected $d_k^I$} \\
\midrule
High (match $> 0.7$) & Large, significant cluster-level improvement & $d > 0.5$ (medium) \\
Moderate (match 0.4-0.7) & Modest cluster improvement & $d = 0.2$--$0.5$ (small-medium) \\
Low (match $< 0.4$) & Minimal or no cluster improvement & $d < 0.2$ (small) \\
Zero match (match $\approx 0$) & Effect = baseline (segmentation adds no value) & $d \approx d_{\text{baseline}}$ \\
\bottomrule
\end{tabular}
\caption{Predicted relationship between intervention-cluster match and effect size}
\label{tab:match-effect-sizes}
\end{table}

\textbf{Test Strategy (Resolving Your Question):}
When comparing segmented vs. non-segmented groups:

\begin{enumerate}
\item \textbf{Non-segmented group:} All clusters receive identical intervention $I_{\text{generic}}$
  \begin{equation}
  E[\Delta O_{\text{non-seg}}] = \frac{1}{K} \sum_k E[\Delta O_k^{I_{\text{generic}}}] = E[\Delta O_{\text{baseline}}]
  \end{equation}

\item \textbf{Segmented group:} Cluster $k$ receives matched intervention $I_k^*$
  \begin{equation}
  E[\Delta O_{\text{seg}}] = \frac{1}{K} \sum_k E[\Delta O_k^{I_k^*}] = E[\Delta O_{\text{baseline}}] + \beta^{\text{match}} \cdot \frac{1}{K}\sum_k \text{Leverage}(I_k^*, C_k)
  \end{equation}

\item \textbf{Predicted difference:}
  \begin{equation}
  E[\Delta O_{\text{seg}}] - E[\Delta O_{\text{non-seg}}] = \beta^{\text{match}} \cdot \text{AvgLeverage}
  \end{equation}

If $\text{AvgLeverage} \approx 0.2$ (weak matching) and $\beta^{\text{match}} = 0.3$ (moderate multiplier effect),
then even perfectly implemented segmentation yields small gains ($\approx 0.06$ effect size).

\end{enumerate}

\textbf{Diagnostic: When Results Are Equal}

If segmented and non-segmented perform equally ($E[\Delta O_{\text{seg}}] \approx E[\Delta O_{\text{non-seg}}]$):

\begin{enumerate}
\item \textbf{Compute Leverage} for each cluster and intervention pair
\item \textbf{If average $\text{Leverage} < 0.3$:} Interventions were poorly matched to clusters
  \begin{itemize}[nosep]
  \item Fix: Re-design interventions to target actual cluster constraints
  \item Segmentation theory (Axioms 1-7) is likely valid
  \end{itemize}
\item \textbf{If average $\text{Leverage} > 0.6$ but outcomes still equal:}
  \begin{itemize}[nosep]
  \item Problem: Segmentation or cluster identification likely flawed (test Axioms 1-7)
  \item Or: Implementation fidelity very poor (interventions didn't match design)
  \end{itemize}
\item \textbf{If $\text{AvgLeverage} = 0$ (all interventions generic):}
  \begin{itemize}[nosep]
  \item No segmented intervention designed $\Rightarrow$ no reason to expect difference
  \item You have created a false control: "segmentation" that produces no differentiated intervention
  \end{itemize}
\end{enumerate}

\textbf{Critical Insight:}
Your observation (``aus der Segmentierung leite ich ja die Massnahmen ab'') is formalized here.
Axiom 9 specifies exactly HOW to derive interventions from segments: match intervention design
to the cluster's leverage profile. Without high leverage matching, segmentation adds no value---
not because segmentation is wrong, but because intervention design is uncoupled from segment characteristics.

\end{tcolorbox}

%%%%%%%%%%%%%%%%%%%%%%%%%%%%%%%%%%%%%%%%%%%%%%%%%%%%%%%%%%%%%%%%%%%%%%%%%%%%%%%
\section{Why 5 Segments in Some Domains: A Special Case}
\label{sec:ba-why-five}
%%%%%%%%%%%%%%%%%%%%%%%%%%%%%%%%%%%%%%%%%%%%%%%%%%%%%%%%%%%%%%%%%%%%%%%%%%%%%%%

In environmental behavior, technology adoption, and health behavior, we often
observe approximately 5 natural clusters. Why?

\subsection{The Rogers Diffusion Mechanism}

When a new behavior diffuses through a population over time, and when adoption
involves both \textit{individual decision-making} and \textit{social influence},
the population naturally stratifies into:

\begin{enumerate}
\item \textbf{Innovators} (~2.5%): Very high $\omega_{\text{novelty}}$, low risk aversion
\item \textbf{Early Adopters} (~13.5%): High social awareness, moderate risk tolerance
\item \textbf{Early Majority} (~34%): Moderate all dimensions (pragmatic)
\item \textbf{Late Majority} (~34%): High social influence sensitivity, moderate innovativeness
\item \textbf{Laggards} (~16%): Very high risk aversion, low innovativeness
\end{enumerate}

This is not a law of nature but a consequence of:
- Normally distributed innovativeness across populations
- Adoption curves that follow S-shaped (logistic) trajectories
- Social reinforcement mechanisms

\subsection{Conditions for 5-Cluster Structure}

The 5-cluster Rogers structure emerges when:
\begin{itemize}[nosep]
\item The behavior is \textit{new or unfamiliar} (high awareness heterogeneity)
\item Adoption involves \textit{individual choice} (not mandatory)
\item \textit{Social influence} is present (people observe peers)
\item The population is \textit{reasonably large} ($N \geq 1000$)
\end{itemize}

When these conditions change, the cluster structure changes:

\begin{table}[h]
\centering
\small
\begin{tabular}{llll}
\toprule
\textbf{Condition Change} & \textbf{Effect} & \textbf{New Cluster Count} & \textbf{Example} \\
\midrule
Mandatory (not voluntary) & Risk aversion becomes less salient & 3-4 & COVID vaccines \\
Deeply identity-threatening & Identity-based polarization & 2-3 & Controversial politics \\
High-cost behavior & Financial constraints split pragmatists & 6-7 & Home solar panels \\
Well-established behavior & Laggards adopt, innovators irrelevant & 2-3 & Smartphone usage \\
\bottomrule
\end{tabular}
\caption{How context changes the natural cluster structure}
\label{tab:cluster-variance}
\end{table}

\textbf{Conclusion:} The 5-cluster Rogers structure is \textit{common but not universal}.
It emerges under specific conditions. Change the conditions, and the cluster
structure changes.

%%%%%%%%%%%%%%%%%%%%%%%%%%%%%%%%%%%%%%%%%%%%%%%%%%%%%%%%%%%%%%%%%%%%%%%%%%%%%%%
\section{Integration with CORE Framework}
\label{sec:ba-integration}
%%%%%%%%%%%%%%%%%%%%%%%%%%%%%%%%%%%%%%%%%%%%%%%%%%%%%%%%%%%%%%%%%%%%%%%%%%%%%%%

Segment emergence connects to all 8 CORE dimensions:

\subsection{WHO (Appendix WHO): Levels and Segment Distributions}

Clustering occurs \textit{within each welfare level}. A firm's employees, a
nation's citizens, and a household's members all have different segment distributions.
The weight function $\alpha^L(\Psi)$ aggregates across levels, and segmentation
operates at each level independently.

\subsection{WHAT (Appendix WAT): Utility Dimensions Drive Clustering}

Segments form along FEPSDE dimensions. Clustering distance metric is typically:
\begin{equation}
d_i \propto \sqrt{\sum_{d \in \{F,E,P,S,D,E\}} (\omega^d_i - \omega^d_{\text{prototype}})^2}
\end{equation}

Different dimensions dominate in different domains (Financial dimension in
economic decisions; Social dimension in peer-effects contexts).
\cite{fehr1999theory, fehr2002cooperation}

\subsection{HOW (Appendix HOW): Complementarities Shape Preference Prototypes}

Cluster prototypes $\mu_k$ are not arbitrary. They reflect natural complementarity
structures: individuals who highly value Financial gains also tend to value
Practical outcomes (complementary motivations). Clustering respects these
complementarity patterns.

\subsection{WHEN (Appendix CTW): Context Determines Cluster Structure}

From Axiom 3: Different contexts produce different clustering. $\Psi$-modulated
heterogeneity directly produces context-specific segments.

\subsection{WHERE (Appendix EST): Parameters Estimated Per Cluster}

All utility parameters, awareness elasticities, and willingness functions can
be estimated per cluster, not just at population level. This increases precision.

\subsection{AWARE (Appendix AWA): Awareness Creates Cluster Visibility}

Without awareness, individuals cannot perceive that alternatives exist. As
$A_i(t)$ increases, individuals become aware of different preference clusters
and can transition between them.

\subsection{READY (Appendix REA): Willingness Operates Within and Across Clusters}

Each cluster has its own willingness threshold $\theta_k$. Transitions between
clusters require both awareness (seeing the alternative) and willingness
(accepting the transition cost).

\subsection{STAGE (Appendix STA): Behavioral Change Journey Differs By Cluster}

The BCJ (5-stage journey) may have different time constants $\tau_k$ for
different clusters. Leaders move faster; Laggards slower.

%------------------------------------------------------------------------------
% HHH (METHOD-TOOLKIT) INTEGRATION - NEW JANUARY 2026
%------------------------------------------------------------------------------
\subsection{HHH (METHOD-TOOLKIT): Segment-Specific Interventions}
\label{sec:ba-hhh-integration}

\begin{tcolorbox}[colback=orange!5!white, colframe=orange!75!black,
    title=\textbf{Integration: Segments $\leftrightarrow$ HHH (Intervention Toolkit)}]

\textbf{Why This Matters:}
Appendix HHH Axiom HHH-A5 (Segment Specificity) formalizes what BA establishes empirically:
the same intervention has dramatically different effects on different segments.

\textbf{Segment Multipliers $\sigma_s$ (from HHH Axiom HHH-A5):}
\begin{equation}
E(\mathcal{I}, s) = E_{\text{base}}(\mathcal{I}) \cdot \sigma_s(\mathcal{I})
\end{equation}
where $\sigma_s \in [-0.5, 2.0]$ --- negative values indicate reactance.

\textbf{Segment-Intervention Effectiveness Matrix:}
\begin{center}
\renewcommand{\arraystretch}{1.2}
\footnotesize
\begin{tabular}{@{}lccccc@{}}
\toprule
\textbf{HHH Type} & \textbf{Leaders} & \textbf{Social} & \textbf{Pragmatic} & \textbf{Autonomy} & \textbf{Skeptics} \\
\midrule
Informative (1) & $\sigma=1.2$ & $\sigma=0.8$ & $\sigma=1.0$ & $\sigma=1.0$ & $\sigma=0.6$ \\
Normative (2) & $\sigma=0.7$ & $\sigma=1.6$ & $\sigma=1.0$ & $\sigma=-0.3$ & $\sigma=0.4$ \\
Structural (3) & $\sigma=0.8$ & $\sigma=1.0$ & $\sigma=1.4$ & $\sigma=0.7$ & $\sigma=0.5$ \\
Incentive (4) & $\sigma=0.6$ & $\sigma=0.7$ & $\sigma=1.5$ & $\sigma=1.2$ & $\sigma=0.8$ \\
Commitment (5) & $\sigma=1.4$ & $\sigma=1.0$ & $\sigma=0.8$ & $\sigma=0.5$ & $\sigma=0.3$ \\
Identity (7) & $\sigma=1.5$ & $\sigma=1.3$ & $\sigma=0.6$ & $\sigma=1.4$ & $\sigma=0.4$ \\
\bottomrule
\end{tabular}
\end{center}

\textbf{Key Insight:} Normative interventions \textbf{backfire} ($\sigma < 0$) on autonomy-seeking segments!

\textbf{What Segments Provide to HHH:}
\begin{itemize}[nosep]
\item Segment definitions for Field F9 specification
\item Segment-specific behavioral parameters ($\beta_s$, $\delta_s$)
\item $\sigma_s$ multipliers for effectiveness prediction
\end{itemize}

\textbf{What HHH Provides to Segments:}
\begin{itemize}[nosep]
\item Segment-targeted intervention portfolios
\item Complementarity-aware combinations per segment
\item Journey-segment interaction for optimal timing
\end{itemize}

\textbf{Interface:} HHH Field F9 (Target Segment) $\leftrightarrow$ BA segment $s \in \mathcal{S}$

\textbf{Cross-Reference:} See Appendix HHH Section 4 (Axiom HHH-A5) for formal specification.

\end{tcolorbox}

%%%%%%%%%%%%%%%%%%%%%%%%%%%%%%%%%%%%%%%%%%%%%%%%%%%%%%%%%%%%%%%%%%%%%%%%%%%%%%%
\section{Worked Example: Environmental Behavior in Three Contexts}
\label{sec:ba-example}
%%%%%%%%%%%%%%%%%%%%%%%%%%%%%%%%%%%%%%%%%%%%%%%%%%%%%%%%%%%%%%%%%%%%%%%%%%%%%%%

This example shows how the same behavior (sustainable consumption) produces
different cluster structures in different contexts.

\subsection{Context 1: Novel Technology (Electric Vehicles)}

\textbf{Situation:} EVs are still new in most markets. High awareness heterogeneity.
Strong social proof effects. Personal cost significant.

\textbf{Clustering Result:} 5 natural clusters emerge.

\begin{itemize}[nosep]
\item \textbf{Innovators} (2-3\%): Early adopters, high income, want social status
\item \textbf{Early Adopters} (12-15\%): Environmentally conscious, willing to pay premium
\item \textbf{Pragmatists} (30-35\%): Will adopt if payback period is 5-7 years
\item \textbf{Late Majority} (30-35\%): Adopt when cost parity with ICE achieved
\item \textbf{Laggards} (15-20\%): Resistant to change, prefer familiar technology
\end{itemize}

\textbf{Segmentation Quality:} Silhouette score $\approx 0.65$ (good). Clusters
differ significantly in income, environmental concern, and risk tolerance.

\subsection{Context 2: Mandatory Behavior (COVID Vaccination)}

\textbf{Situation:} Vaccination is strongly recommended (or mandatory in some
regions). High stakes (health). Politicized.

\textbf{Clustering Result:} Only 2-3 clusters emerge.

\begin{itemize}[nosep]
\item \textbf{Compliers} (70-75\%): Follow expert guidance, want to be safe
\item \textbf{Skeptics} (15-20\%): Distrust authorities, prefer natural immunity
\item (Optional) \textbf{Defiant} (5-10\%): Actively oppose vaccination on principle
\end{itemize}

\textbf{Segmentation Quality:} Silhouette score $\approx 0.55$ (moderate). Risk
aversion dominates; innovativeness irrelevant. Fewer natural clusters.

\subsection{Context 3: Established Behavior (Smartphone Usage)}

\textbf{Situation:} Smartphones are ubiquitous and affordable. Adoption is no
longer a choice but a necessity. Segmentation now focuses on \textit{usage intensity}
and \textit{app preferences}, not adoption.

\textbf{Clustering Result:} 2-3 clusters emerge.

\begin{itemize}[nosep]
\item \textbf{Heavy Users} (30-40\%): Use for work, entertainment, socializing
\item \textbf{Moderate Users} (40-50\%): Use for essential tasks only
\item \textbf{Light Users} (10-20\%): Use only for calls and texts
\end{itemize}

\textbf{Segmentation Quality:} Silhouette score $\approx 0.60$ (good, but lower
than EV case). Segmentation now is based on usage intensity, not adoption
decision.

\subsection{Key Insight}

The same population can have:
- 5 clusters for novel, voluntary, costly behaviors
- 3 clusters for mandatory, high-stakes behaviors
- 2-3 clusters for established, ubiquitous behaviors

This is **not because the underlying theory changed**, but because **context
changed the heterogeneity distribution** (Axiom 3).

%%%%%%%%%%%%%%%%%%%%%%%%%%%%%%%%%%%%%%%%%%%%%%%%%%%%%%%%%%%%%%%%%%%%%%%%%%%%%%%
\section{Implications for Practice}
\label{sec:ba-practice}
%%%%%%%%%%%%%%%%%%%%%%%%%%%%%%%%%%%%%%%%%%%%%%%%%%%%%%%%%%%%%%%%%%%%%%%%%%%%%%%

\begin{tcolorbox}[colback=green!5!white, colframe=green!75!black,
    title=Practical Implications of Segmentation Genesis]

\begin{enumerate}

\item \textbf{Never assume a fixed number of segments.} Conduct data-driven
clustering analysis for your specific domain and context. Use elbow method,
silhouette scores, or BIC to determine optimal $k$.

\item \textbf{Segments are emergent, not theoretical constructs.} Don't impose
preconceived categories (``these 5 are the true segments''). Let the data reveal
the natural structure.

\item \textbf{Context is everything.} The same behavior in different contexts
produces different segments. Electric vehicles in wealthy vs. developing nations
show different clustering structures. Always analyze context-specific segmentation.

\item \textbf{Use standard clustering algorithms.} K-means, hierarchical clustering,
and spectral clustering are validated, widely available, and produce stable results.
Don't invent custom segmentation schemes.

\item \textbf{Validate cluster differences in intervention response.} A good
segmentation scheme must demonstrate that clusters respond differently to
interventions. If all clusters respond identically, segmentation adds no value.

\item \textbf{Monitor cluster transitions.} Track not just overall adoption but
inter-cluster flows. Which clusters are gaining/losing members? Why? This reveals
whether interventions are triggering transitions.

\item \textbf{Design segment-specific interventions only after clustering.} Map
intervention levers to cluster characteristics. What works for cluster A may be
useless or counterproductive for cluster B.
\cite{thaler2008nudge, dellavigna2021illusion}

\end{enumerate}

\end{tcolorbox}

%%%%%%%%%%%%%%%%%%%%%%%%%%%%%%%%%%%%%%%%%%%%%%%%%%%%%%%%%%%%%%%%%%%%%%%%%%%%%%%
\section{Results: Consequences of the Axioms}
\label{sec:ba-results}
%%%%%%%%%%%%%%%%%%%%%%%%%%%%%%%%%%%%%%%%%%%%%%%%%%%%%%%%%%%%%%%%%%%%%%%%%%%%%%%

The six axioms of segment emergence yield powerful conclusions for understanding
behavioral heterogeneity:

\subsection{Result 1: Universal Applicability}
Axioms 1-6 apply to every behavior-change domain, every population, and every
context. There are no exceptions. This means segmentation is always possible and
always necessary.

\subsection{Result 2: Context-Dependence}
While the axioms are universal, their outcomes are not. The same axioms applied
to different populations/contexts produce different cluster counts ($k$), cluster
sizes, cluster prototypes, and cluster transition dynamics.

\subsection{Result 3: Data-Driven Necessity}
Because $k$ is context-dependent and unknowable a priori, segmentation MUST be
empirical. Imposing a fixed number of segments (e.g., ``always 5'') violates the
axioms and produces poor segmentation quality.

\subsection{Result 4: Empirical Validity}
A segmentation scheme is valid only if clusters demonstrate different intervention
response patterns (Axiom 4). Mathematical cluster quality alone is insufficient;
behavioral homogeneity must be empirically verified.

%%%%%%%%%%%%%%%%%%%%%%%%%%%%%%%%%%%%%%%%%%%%%%%%%%%%%%%%%%%%%%%%%%%%%%%%%%%%%%%
\section{Key Propositions}
\label{sec:ba-propositions}
%%%%%%%%%%%%%%%%%%%%%%%%%%%%%%%%%%%%%%%%%%%%%%%%%%%%%%%%%%%%%%%%%%%%%%%%%%%%%%%

\begin{proposition}[Universality of Heterogeneity]
For any population of size $N > 100$ and any behavior-change domain, heterogeneity
in utility preferences, awareness, and willingness exceeds zero-noise threshold.
Therefore, Axiom 1 applies universally.
\end{proposition}

\begin{proposition}[Cluster Detectability]
If Axiom 1 holds and population is large enough ($N > 500$), standard clustering
algorithms will reliably detect $k \geq 2$ natural clusters with high probability
(statistical power $>0.8$).
\end{proposition}

\begin{proposition}[Context Determines $k$]
The optimal number of clusters $k^*$ is a function of context: $k^* = f(\Psi)$.
Two populations with identical preference distributions but different contexts
will show different optimal $k$.
\end{proposition}

\begin{proposition}[Cluster Homogeneity Implies Intervention Efficiency]
If Axiom 4 holds (within-cluster behavioral homogeneity), then segment-specific
interventions yield higher expected utility gains than population-average
interventions by a factor of $1 + \rho$, where $\rho > 0$ is the
within-cluster correlation in intervention response.
\end{proposition}

\begin{proposition}[Rogers Structure as Special Case]
The 5-cluster Rogers structure (Innovators, Early Adopters, Early Majority,
Late Majority, Laggards) emerges when: (a) behavior is voluntary, (b) adoption
involves social learning, and (c) adoption costs are moderate-to-high relative
to income. Remove any condition, and $k \neq 5$.
\end{proposition}

\begin{proposition}[Cluster Size Regularity under Multivariate Normality]
If heterogeneity in preference space is approximately multivariate normal
$\mathcal{H} \sim \mathcal{N}(\mu_{\mathcal{H}}, \Sigma_{\mathcal{H}})$, then cluster sizes
$|C_k|$ exhibit predictable statistical properties:
\begin{enumerate}[nosep]
\item Mean cluster size: $E[|C_k|] = N/k^*$
\item Coefficient of variation: $CV(|C_k|) \in [0.3, 0.5]$
\item Distribution shape: approximately Binomial or Beta
\item Probability that largest cluster exceeds $N/k^* + 2\sigma$: $< 5\%$ (under normality)
\end{enumerate}
Deviations from this pattern (e.g., one cluster $> 60\%$ of population) indicate either
non-normal heterogeneity distribution (bimodal, heavily skewed) or suboptimal clustering.
\end{proposition}

\begin{proposition}[Rogers Curve as Mathematical Consequence of Normal Heterogeneity + Social Learning]
The 5-cluster Rogers distribution $\{2.5\%, 13.5\%, 34\%, 34\%, 16\%\}$ emerges from:
\begin{enumerate}[nosep]
\item \textbf{Input:} Multivariate normal heterogeneity in innovativeness $\omega_{\text{nov}}$
\item \textbf{Process:} S-curve adoption dynamics with social influence proportional to adoption prevalence
\item \textbf{Mechanism:} Each individual's adoption time is a function of $\omega_{\text{nov}}(i)$ and cumulative adoption at time $t$
\item \textbf{Result:} Natural stratification produces Beta distribution with $\alpha \approx 2, \beta \approx 3$
\end{enumerate}
This is \textit{not} a law of nature but a \textit{mathematical consequence} of normal heterogeneity
coupled with sigmoid diffusion dynamics. Change the heterogeneity distribution shape, and the
cluster size distribution changes accordingly.
\cite{rogers2003diffusion, montoya2003extensions}
\end{proposition}

%----------- AXIOM 10 -----------
\begin{tcolorbox}[colback=purple!5!white, colframe=purple!75!black,
    title=Axiom SEG-VAL-10: Theory Externalization Principle]

\textbf{Statement:}
Experimental design forces practitioners to convert implicit behavioral theories into explicit mathematical
structures. By requiring specification of leverage profiles and intervention design profiles before data collection,
experiments \textit{externalize} tacit knowledge and make it falsifiable. The mechanism is universal: practitioners
cannot design an experiment without articulating the exact mechanisms they believe are operating.
\cite{popper1959logic, angrist2010mostly}

\textbf{Formal Definition:}
Before any experiment comparing segmented vs.\ non-segmented interventions, the practitioner must specify
two mathematical objects:

\begin{equation}
\text{(a) Leverage Profile:} \quad \mathcal{L}_k = (\ell_k^d)_{d \in \{F,E,P,S,D,E\}}
\quad \text{where} \quad \ell_k^d = \frac{\text{Var}(\omega_i^d | i \in C_k)}{\sum_{d'} \text{Var}(\omega_i^{d'} | i \in C_k)}
\end{equation}

\begin{equation}
\text{(b) Intervention Design Profile:} \quad \mathcal{D}_I = (\delta_I^d)_{d \in \{F,E,P,S,D,E\}} \quad \text{where} \quad \delta_I^d \in [0,1]
\end{equation}

This requirement forces the implicit intuition (e.g., ``identity-driven clusters respond to identity appeals'')
to become explicit: Which dimensions actually matter for this cluster? ($\ell_k^d$ specified). How intensely does
the intervention target each dimension? ($\delta_I^d$ specified).

\textbf{Interpretation:}
Without Axiom 10, practitioners remain in the realm of craft knowledge: tacit, non-transferable, undocumented.
Axiom 10 creates a forcing function: to design an experiment, you \textit{must} externalize your theory.
This externaliza

tion is the bridge from implicit to explicit knowledge.

\textbf{Consequence:}
Practitioners cannot hide behind vague intuitions when designing experiments. They must commit to specific
mathematical relationships between cluster characteristics and intervention effectiveness. This commitment
enables the subsequent axioms (11, 12, 13) to function.

\end{tcolorbox}

\vspace{1em}

%----------- AXIOM 11 -----------
\begin{tcolorbox}[colback=green!5!white, colframe=green!75!black,
    title=Axiom SEG-VAL-11: Pre-Registration and Falsifiability Principle]

\textbf{Statement:}
Pre-registration transforms externalized theories into falsifiable hypotheses. By committing in advance to
specific predictions, effect sizes, and diagnostic procedures, practitioners create an ``asymmetry'' that
enables genuine scientific testing. Results either confirm or disconfirm the theory; they cannot be ambiguous
post-hoc rationalizations.
\cite{nosek2018registered, open2015estimating}

\textbf{Formal Definition:}
Before data collection, the practitioner pre-registers:

\begin{enumerate}
\item \textbf{Primary Prediction:}
\begin{equation}
E[\Delta O_{\text{seg}}] - E[\Delta O_{\text{non-seg}}] = \beta^{\text{match}} \cdot \text{AvgLeverage}
\end{equation}

\item \textbf{Expected $\beta^{\text{match}}$:} Quantified from literature or prior pilots (e.g., $\beta^{\text{match}} = 0.15$)

\item \textbf{Expected AvgLeverage:} Computed from (a) and (b) in Axiom 10

\item \textbf{Diagnostic Protocol:} Pre-specified decision tree:
\begin{itemize}[nosep]
\item IF $\text{AvgLeverage}_{\text{observed}} < 0.3$ AND outcomes equal $\Rightarrow$ Intervention matching problem
\item IF $\text{AvgLeverage}_{\text{observed}} > 0.6$ AND outcomes equal $\Rightarrow$ Segmentation or fidelity problem
\end{itemize}

\end{enumerate}

\textbf{Interpretation:}
Pre-registration transforms ``what do I think will happen?'' into ``what exactly must happen for my theory to be correct?''
This creates falsifiability. The theory is no longer infinitely flexible; it makes specific, testable predictions.

\textbf{Critical Implication:}
Without pre-registration (Axiom 11), the externalized theory (Axiom 10) remains only partially constrained.
Practitioners could still rationalize any outcome post-hoc. Axiom 11 prevents this by locking in predictions
before seeing results.

\end{tcolorbox}

\vspace{1em}

%----------- AXIOM 12 -----------
\begin{tcolorbox}[colback=orange!5!white, colframe=orange!75!black,
    title=Axiom SEG-VAL-12: Diagnostic Refinement Principle]

\textbf{Statement:}
Experimental results provide more than binary validation (true/false). They provide \textbf{diagnostic signals} that
identify exactly WHERE theory failed and HOW parameters should be refined. By computing actual leverage matches
and comparing them to predictions, practitioners can distinguish between segmentation errors, intervention design
errors, and implementation fidelity errors. This enables cumulative theory refinement rather than wholesale rejection.
\cite{button2013power, bland1995multiple, camerer2018evaluating}

\textbf{Formal Definition:}
If $E[\Delta O_{\text{seg}}] \approx E[\Delta O_{\text{non-seg}}]$ (predicted difference not observed):

\begin{enumerate}

\item \textbf{Compute Diagnostic Signal:}
\begin{equation}
\text{AvgLeverage}_{\text{observed}} = \frac{1}{K} \sum_k \text{Leverage}(I_k^*, C_k) \quad \text{from empirical data}
\end{equation}

\item \textbf{Compare to Prediction:}
\begin{equation}
\Delta \text{Leverage} = \text{AvgLeverage}_{\text{observed}} - \text{AvgLeverage}_{\text{predicted}}
\end{equation}

\item \textbf{Diagnosis from Diagnostic Signal:}
\begin{itemize}[nosep]
\item $\text{AvgLeverage}_{\text{observed}} < 0.3$: Intervention design was poorly matched to clusters
  (Solution: redesign $\delta_I^d$, keep segmentation and Axioms 1-7)
\item $0.4 < \text{AvgLeverage}_{\text{observed}} < 0.6$ AND low effect: Interventions partly matched
  (Solution: increase specificity, recalibrate $\delta_I^d$ for higher-leverage dimensions)
\item $\text{AvgLeverage}_{\text{observed}} > 0.6$ AND low effect: Segmentation or implementation fidelity failed
  (Solution: revalidate clusters against Axioms 1-7, or improve treatment fidelity)
\end{itemize}

\end{enumerate}

\textbf{Interpretation:}
Axiom 12 operationalizes the difference between ``my theory is wrong'' vs.\ ``my theory is incompletely specified.''
If actual leverage is low but predicted was high, the problem is parameter estimation (recalibrate $\ell_k^d$ or $\delta_I^d$),
not theoretical structure. If actual leverage matches prediction but outcomes still don't change, the problem is elsewhere
(segmentation quality or implementation).

\textbf{Consequence:}
This diagnostic structure enables \textbf{targeted} refinement instead of wholesale theory rejection. Practitioners
learn exactly which parameters need adjustment and can refine them for the next experiment cycle.

\end{tcolorbox}

\vspace{1em}

%----------- AXIOM 13 -----------
\begin{tcolorbox}[colback=cyan!5!white, colframe=cyan!75!black,
    title=Axiom SEG-VAL-13: Cumulative Knowledge Principle]

\textbf{Statement:}
Externalized, falsified, and refined theories are \textbf{transferable} and \textbf{generalizable} across practitioners
and contexts. The cycle of externalization (Axiom 10) → falsification (Axiom 11) → refinement (Axiom 12) creates
cumulative scientific knowledge. Unlike implicit craft knowledge (which is lost when experts retire), explicit theories
accumulate, transfer, and improve as each experiment cycle adds data and constraints.
\cite{camerer2018evaluating, open2015estimating, nosek2015power}

\textbf{Formal Definition:}
After completing an experiment cycle, the refined theory can be:

\begin{enumerate}

\item \textbf{Transferred to Other Practitioners:}
Operational specification: ``In domain X, clusters of type C respond to interventions with $(\delta_I^{d*}, \beta_I^{b*})$ to achieve
average leverage of L$^*$, producing effect size d$^*$.'' This is teachable because all parameters are explicit.

\item \textbf{Tested in New Contexts:}
Same measurements ($\ell_k^d, \delta_I^d$), same predictions ($\beta^{\text{match}}, E[\Delta O]$), different population/context.
Enables meta-analysis across contexts.

\item \textbf{Refined with Each New Data Point:}
Bayesian updating of $\beta^{\text{match}}$ as more experiments accumulate. Prior is refined theory from previous cycle;
likelihood is new empirical data; posterior becomes new refined theory.

\item \textbf{Accumulated in Registry (FFF):}
Seed models stored with: (a) cluster leverage profiles, (b) intervention design profiles, (c) empirical $\beta^{\text{match}}$ estimates,
(d) success/failure patterns by context.

\end{enumerate}

\textbf{Interpretation:}
Axiom 13 is the bridge from individual practitioner knowledge to \textbf{cumulative scientific knowledge}.
It explains why the externalization-falsification-refinement cycle (Axioms 10-12) produces science rather than just applied experimentation.
Science accumulates; craft does not.

\textbf{Comparison to Implicit Theory:}

\begin{table}[h]
\centering
\small
\begin{tabular}{lll}
\toprule
\textbf{Property} & \textbf{Implicit Theory} & \textbf{Explicit Theory (Axioms 10-13)} \\
\midrule
Transferable? & No (expert-dependent) & Yes (operationally specified) \\
Cumulative? & No (lost when expert leaves) & Yes (registry accumulates) \\
Falsifiable? & No (always rationalizable) & Yes (pre-registered predictions) \\
Refinable? & No (stuck in original form) & Yes (diagnostic refinement) \\
Scientific? & Craft & Science \\
\bottomrule
\end{tabular}
\caption{Why Axioms 10-13 create science from craft knowledge}
\label{tab:implicit-vs-explicit}
\end{table}

\textbf{Final Insight:}
Your observation (``aber natürlich auch explizite Theory ständig zu verbessern'') is formalized in Axioms 10-13.
Experiments don't just validate theories; they \textit{externalize}, \textit{falsify}, and \textit{refine} them continuously.
Each experiment cycle produces knowledge that transfers and accumulates.

\end{tcolorbox}

\vspace{1.5em}

%%%%%%%%%%%%%%%%%%%%%%%%%%%%%%%%%%%%%%%%%%%%%%%%%%%%%%%%%%%%%%%%%%%%%%%%%%%%%%%
\section{Integration: Axioms 1-13 as a Complete Segmentation Science Framework}
\label{sec:ba-complete-framework}
%%%%%%%%%%%%%%%%%%%%%%%%%%%%%%%%%%%%%%%%%%%%%%%%%%%%%%%%%%%%%%%%%%%%%%%%%%%%%%%

The thirteen axioms form a hierarchical, complete framework for segmentation science:

\begin{tcolorbox}[colback=cyan!10!white, colframe=cyan!75!black,
    title=13-Axiom Hierarchical Framework]

\textbf{Layer 1: Segment Emergence (Axioms 1-7)}
\begin{itemize}[nosep]
\item \textbf{WHY} segments exist: Heterogeneity is universal (Axiom 1), produces natural clustering (Axiom 2), with cluster sizes following predictable distributions (Axiom 7)
\item \textbf{WHERE} cluster count is determined: By context Ψ (Axiom 3), with cluster homogeneity required (Axiom 4)
\item \textbf{HOW} to identify them: Via clustering algorithms (Axiom 6), with transitions possible (Axiom 5)
\end{itemize}

\textbf{Layer 2: Intervention Matching (Axioms 8-9)}
\begin{itemize}[nosep]
\item \textbf{WHAT} are leverage points: Each cluster has identifiable leverage profiles $\ell_k^d$ (Axiom 8)
\item \textbf{HOW} to design interventions: Match $\delta_I^d$ to cluster leverage profiles (Axiom 8)
\item \textbf{WHEN} segmentation adds value: When $\text{Leverage}(I, C_k)$ is high (Axiom 9)
\item \textbf{DIAGNOSTIC}: When results equal, compute $\text{AvgLeverage}$ to identify problem source
\end{itemize}

\textbf{Layer 3: Experimental Validation (Axioms 10-13)}
\begin{itemize}[nosep]
\item \textbf{HOW} to externalize implicit theory: Specify $\mathcal{L}_k$ and $\mathcal{D}_I$ mathematically (Axiom 10)
\item \textbf{MAKE} theory testable: Pre-register predictions and diagnostic trees (Axiom 11)
\item \textbf{REFINE} parameters: Use diagnostic signals ($\text{AvgLeverage}_{\text{observed}}$) to identify and fix errors (Axiom 12)
\item \textbf{ACCUMULATE} knowledge: Transfer refined theories across practitioners and contexts (Axiom 13)
\end{itemize}

\textbf{Net Result:}
From implicit intuitions to explicit, falsifiable, transferable, and continuously improving scientific knowledge.
This is why practitioners cannot hide behind vagueness when using this framework.

\end{tcolorbox}

\vspace{1em}

%%%%%%%%%%%%%%%%%%%%%%%%%%%%%%%%%%%%%%%%%%%%%%%%%%%%%%%%%%%%%%%%%%%%%%%%%%%%%%%
\section{Summary}
\label{sec:ba-summary}
%%%%%%%%%%%%%%%%%%%%%%%%%%%%%%%%%%%%%%%%%%%%%%%%%%%%%%%%%%%%%%%%%%%%%%%%%%%%%%%

\begin{tcolorbox}[colback=green!10!white, colframe=green!75!black,
    title=\textbf{Appendix PAP1: Summary}]

\textbf{Core Question:} How do behavioral segments emerge from heterogeneous populations?

\textbf{Main Contribution:}
Segments are not theoretical constructs imposed by researchers. They are
\textit{natural clusters} that emerge from heterogeneity in utility preferences (ω),
awareness (A), willingness (W), and context (Ψ). Thirteen universal axioms govern
this process: Axioms 1-7 formalize segment \textit{emergence}, Axioms 8-9 formalize how to
\textit{match interventions} to leverage segment differences, and Axioms 10-13 formalize how
\textit{experiments externalize, validate, refine, and accumulate} segmentation science.
The framework is now complete: it operationalizes the entire journey from implicit craft knowledge
to explicit, falsifiable, transferable, and continuously improving science.

\textbf{Key Insight:}
The number and type of segments are NOT fixed universally. Instead, they
emerge data-driven from each specific population/context combination. Some
domains naturally produce 3-4 clusters; others 6-7. The famous ``5 segments''
(Rogers diffusion curve) are a special case, not a universal law.

\textbf{Universal Principles (Context-Independent):}
\begin{itemize}[nosep]
\item \textbf{Emergence:} Heterogeneity universal (Axiom 1) → natural clustering (Axiom 2) with context-dependent count (Axiom 3), cluster homogeneity (Axiom 4), transitions (Axiom 5), empirical identifiability (Axiom 6), predictable sizes (Axiom 7)
\item \textbf{Matching:} Leverage profiles identifiable (Axiom 8) → effect sizes scale with match specificity (Axiom 9)
\item \textbf{Validation:} Experiments externalize theories (Axiom 10) → pre-registration falsifies (Axiom 11) → diagnostics refine (Axiom 12) → explicit theories accumulate (Axiom 13)
\end{itemize}

\textbf{Context-Dependent Outcomes:}
\begin{itemize}[nosep]
\item Cluster count: 2-8 (typical range 3-5)
\item Cluster prototypes: Domain-specific
\item Cluster sizes: Distributed with CV ≈ 0.3-0.5 under normal heterogeneity
\item Cluster size patterns: Diagnostic of heterogeneity distribution shape
\item Transition rates: Context-dependent
\end{itemize}

\textbf{Practical Process:}
1. Collect heterogeneity data ($\omega_i, A_i, W_i$ vectors)
2. Determine optimal $k$ via clustering metrics
3. Identify clusters via standard algorithms
4. Validate clusters show different intervention response
5. Compute leverage profile for each cluster (Axiom 8)
6. Design segment-specific interventions matched to leverage profiles
7. Test segmented vs. non-segmented groups with leverage diagnostics (Axiom 9)
8. If equal outcomes, compute $\text{AvgLeverage}$ to diagnose problem (poor matching vs. bad segmentation)

\textbf{The 13-Axiom Framework Completes the Circle:}
\begin{itemize}[nosep]
\item \textbf{Axioms 1-7 answer:} WHERE do segments come from? (population heterogeneity)
\item \textbf{Axioms 8-9 answer:} HOW to exploit segment differences? (matching interventions to leverage profiles)
\item \textbf{Axioms 10-13 answer:} HOW to ensure segmentation science is cumulative, not craft? (externalization → falsification → refinement → accumulation)
\end{itemize}

Without Axioms 10-13, practitioners remain in implicit theory: tacit, non-transferable, non-cumulative.
With Axioms 10-13, practitioners produce explicit theory: operationalized, falsifiable, teachable, cumulative.
This is why your observation was critical: ``segmentation without explicit theory is sinnlos.'' Axioms 10-13
operationalize how experiments transform segmentation from craft into science.

\textbf{Integration:} See Section \ref{sec:ba-complete-framework} for the hierarchical integration of all 13 axioms
as a complete segmentation science framework.  Segments emerge from CORE foundations (utility structure,
awareness, willingness, context) and operationalize heterogeneity in behavior change. Axioms 8-13 enable
empirical validation, continuous theory improvement, and cumulative scientific knowledge.

\end{tcolorbox}

%%%%%%%%%%%%%%%%%%%%%%%%%%%%%%%%%%%%%%%%%%%%%%%%%%%%%%%%%%%%%%%%%%%%%%%%%%%%%%%
\section{Glossary of Symbols}
\label{sec:ba-glossary}
%%%%%%%%%%%%%%%%%%%%%%%%%%%%%%%%%%%%%%%%%%%%%%%%%%%%%%%%%%%%%%%%%%%%%%%%%%%%%%%

\begin{tcolorbox}[colback=blue!5!white, colframe=blue!75!black]
\textbf{Central Reference:} For complete symbol definitions, see
\textbf{Appendix GLS} (\textit{Glossary and Notation}, \S\ref{app:glossary}).

This section lists symbols introduced or specialized in this appendix.
\end{tcolorbox}

\vspace{0.5em}

\begin{table}[htbp]
\centering
\caption{Symbols in Appendix PAP1 (Segmentation Genesis)}
\label{tab:ba-symbols}
\renewcommand{\arraystretch}{1.4}
\begin{tabular}{@{}clp{5cm}c@{}}
\toprule
\textbf{Symbol} & \textbf{Name} & \textbf{Definition} & \textbf{In G?} \\
\midrule
$\mathcal{H}$ & Heterogeneity space & Space of all (ω, A, W) vectors & NEW \\
$C_k$ & Cluster $k$ & Set of individuals in cluster $k$ & NEW \\
$\mu_k$ & Cluster prototype & Center/prototype of cluster $k$ & NEW \\
$k$ & Cluster count & Number of natural clusters (data-driven) & NEW \\
$d(\cdot, \cdot)$ & Distance metric & Metric in heterogeneity space & NEW \\
$\omega_i^d$ & Utility weight & Individual $i$'s weight on dimension $d$ & ✓ \\
$A_i(t)$ & Awareness trajectory & Individual $i$'s awareness over time & ✓ \\
$W_i(t)$ & Willingness function & Individual $i$'s willingness at time $t$ & ✓ \\
$\epsilon_i^I$ & Intervention elasticity & Willingness elasticity to intervention $I$ & NEW \\
$P(C_i \to C_j)$ & Transition probability & Probability of transition from cluster $i$ to $j$ & NEW \\
\bottomrule
\end{tabular}
\end{table}

%%%%%%%%%%%%%%%%%%%%%%%%%%%%%%%%%%%%%%%%%%%%%%%%%%%%%%%%%%%%%%%%%%%%%%%%%%%%%%%
\section{Critical Foundations and Objections}
\label{sec:ba-objections}
%%%%%%%%%%%%%%%%%%%%%%%%%%%%%%%%%%%%%%%%%%%%%%%%%%%%%%%%%%%%%%%%%%%%%%%%%%%%%%%

\subsection{Objection 1: ``Are 5 Segments Not Universal?''}

\textbf{Objection:} Rogers diffusion curve shows 5 segments. Isn't this universal?

\textbf{Response:} Rogers structure is highly common but not universal. It emerges
under specific conditions (voluntary adoption, social learning, moderate cost).
Change any condition (make it mandatory, high-cost, identity-threatening), and
cluster structure changes. Empirical segmentation analysis in each domain reveals
the actual structure.

\subsection{Objection 2: ``Clustering is Data-Driven, Not Theory-Based.'''}

\textbf{Objection:} If clusters are purely empirical, aren't they arbitrary?

\textbf{Response:} The axioms are theoretically grounded. Axioms 1-6 explain
\textit{why} clusters emerge. Clustering algorithms are the \textit{technical}
tool for finding them. This is not arbitrary; it's systematic application of
theory-justified methods.

\subsection{Objection 3: ``How Do We Know We Found the Right $k$?''}

\textbf{Objection:} Clustering metrics (silhouette, BIC) can be ambiguous. How
do we pick the true $k$?

\textbf{Response:} Validation through intervention response. If two cluster
solutions both have high silhouette scores, we compare their intervention
elasticity variance. The solution that better separates intervention responses
is superior. Ultimately, practical utility (not mathematical aesthetics) validates
segmentation.

%%%%%%%%%%%%%%%%%%%%%%%%%%%%%%%%%%%%%%%%%%%%%%%%%%%%%%%%%%%%%%%%%%%%%%%%%%%%%%%
\section{Open Issues and Research Agenda}
\label{sec:ba-open}
%%%%%%%%%%%%%%%%%%%%%%%%%%%%%%%%%%%%%%%%%%%%%%%%%%%%%%%%%%%%%%%%%%%%%%%%%%%%%%%

\subsection{Methodological Priorities}

\begin{enumerate}

\item \textbf{Cluster Stability Across Time:} How stable are clusters as context
evolves? Can we predict cluster migrations before they occur?

\item \textbf{Optimal Distance Metrics:} Which distance metrics in preference
space best predict intervention response? Euclidean? Mahalanobis? Custom metrics?

\item \textbf{Multi-Domain Clustering:} Some individuals may cluster differently
in different domains. How do we handle cross-domain heterogeneity?

\end{enumerate}

\subsection{Empirical Priorities}

\begin{enumerate}

\item \textbf{Cluster Validation RCTs:} Design experiments that randomly assign
individuals within a cluster to different interventions. Do clusters show
predicted response differences?

\item \textbf{Predictive Clustering:} Can we predict cluster membership from
observable characteristics (age, income, past behavior) without collecting full
preference vectors?

\item \textbf{Cross-Cultural Cluster Structures:} Do the same behaviors produce
similar cluster structures across cultures, or are they culture-specific?

\end{enumerate}

%%%%%%%%%%%%%%%%%%%%%%%%%%%%%%%%%%%%%%%%%%%%%%%%%%%%%%%%%%%%%%%%%%%%%%%%%%%%%%%
\section{References}
\label{sec:ba-references}
%%%%%%%%%%%%%%%%%%%%%%%%%%%%%%%%%%%%%%%%%%%%%%%%%%%%%%%%%%%%%%%%%%%%%%%%%%%%%%%

\begin{tcolorbox}[colback=blue!5!white, colframe=blue!75!black]

\textbf{Master Bibliography Reference:}
All citations in this appendix appear in:
\begin{center}
\texttt{bibliography/bcm2\_master\_references.bib}
\end{center}

Key foundational works cited below reference this master bibliography.

\end{tcolorbox}

\vspace{1em}

\nocite{bcm2_master}

\textbf{Clustering and Segmentation Methods:}
\begin{itemize}
\item MacQueen, J. (1967). ``Some methods for classification and analysis of
  multivariate observations.'' \textit{Proceedings of the Fifth Berkeley Symposium}.
  --- Foundational k-means clustering
\item Ward, H.D. (1963). ``Hierarchical clustering.'' \textit{Journal of the
  American Statistical Association}. --- Hierarchical clustering methods
\item Rousseeuw, P.J. (1987). ``Silhouettes: A graphical aid to the interpretation
  and validation of cluster analysis.'' \textit{Journal of Computational and
  Applied Mathematics}. --- Silhouette score validation
\end{itemize}

\textbf{Diffusion of Innovations and Rogers Curve:}
\begin{itemize}
\item Rogers, E.M. (2003). \textit{Diffusion of Innovations} (5th ed.). Free Press.
  --- Classic source on adoption curves and 5-segment structure
\item Bass, F.M. (1969). ``A new product growth model for consumer durables.''
  \textit{Management Science}. --- Bass diffusion model mathematical foundations
\end{itemize}

\textbf{Behavioral Heterogeneity:}
\begin{itemize}
\item Thaler, R.H., \& Sunstein, C.R. (2008). \textit{Nudge}. Penguin.
  --- Recognition of heterogeneous decision-makers
\item Camerer, C.F., Loewenstein, G., \& Rabin, M. (Eds.). (2004).
  \textit{Advances in Behavioral Decision Research}. --- Comprehensive treatment
  of preference heterogeneity
\end{itemize}

%%%%%%%%%%%%%%%%%%%%%%%%%%%%%%%%%%%%%%%%%%%%%%%%%%%%%%%%%%%%%%%%%%%%%%%%%%%%%%%
% CROSS-REFERENCE FOOTER
%%%%%%%%%%%%%%%%%%%%%%%%%%%%%%%%%%%%%%%%%%%%%%%%%%%%%%%%%%%%%%%%%%%%%%%%%%%%%%%

\begin{tcolorbox}[colback=purple!5!white, colframe=purple!75!black,
    title=Cross-References]

\textbf{This Appendix PAP1 references:}
\begin{itemize}[nosep]
\item \textbf{Appendix STA} (CORE-STAGE): BCJ dynamics differ by cluster
\item \textbf{Appendix REA} (CORE-READY): Willingness thresholds differ by cluster
\item \textbf{Appendix AWA} (CORE-AWARE): Awareness distributions shape clustering
\item \textbf{Appendix WAT} (CORE-WHAT): Utility dimensions define clustering space
\item \textbf{Appendix CTW} (CORE-WHEN): Context modulates cluster structure
\item \textbf{Appendix HHH} (REF-SEGMENTATION-HEURISTICS): Quick decision guides for practitioners
\item \textbf{Chapter 14}: Behavioral Change Segments (practical operationalization)
\end{itemize}

\textbf{Referenced by:}
\begin{itemize}[nosep]
\item \textbf{Chapter 14}: For theoretical justification of data-driven clustering
\item \textbf{Chapter 15}: For integration of clusters into intervention design
\item \textbf{Appendix HHH}: For theoretical foundations of diagnostic heuristics
\item \textbf{Appendices EEE, FFF, GGG}: For cluster-aware model design
\end{itemize}

\end{tcolorbox}

\end{document}
